\documentclass[../DoAn.tex]{subfiles}
\begin{document}

Phụ lục này trình bày các khái niệm nền tảng được nhắc tới trong phần thân. Đây là kiến thức chung, không phải đóng góp của đồ án; các khái niệm này được trình bày ở phụ lục để phần thân tóm tắt ngắn gọn và trỏ về các mục dưới.

\section{Blockchain và hợp đồng thông minh}
\label{apdx:blockchain}

Blockchain là một sổ cái phân tán, ghi lại các giao dịch theo từng khối (block) nối tiếp nhau và được nhiều máy cùng giữ bản sao, khiến dữ liệu đã ghi thực tế không thể thay đổi. Ethereum là một blockchain~\cite{wood2014ethereum} cho phép chạy \emph{hợp đồng thông minh} (smart contract) - những đoạn mã được triển khai lên chuỗi và tự động thực thi theo đúng logic đã viết, không ai can thiệp được. Có hai loại tài khoản: tài khoản do người dùng kiểm soát bằng khóa riêng, và tài khoản hợp đồng. Mỗi giao dịch làm thay đổi trạng thái chuỗi đều tốn một khoản phí gọi là phí gas. Token ERC-20 là một chuẩn chung để tạo các đồng tiền số trên Ethereum, quy định các hàm cơ bản như chuyển tiền và cấp quyền chi tiêu.

\section{Nonce}
\label{apdx:nonce}

Nonce là một con số dùng một lần, gắn với mỗi lệnh hoặc mỗi giao dịch, nhằm chống việc phát lại (replay). Nếu không có nonce, một lệnh đã ký có thể bị gửi đi nhiều lần để thực hiện lặp lại ngoài ý muốn của người ký. Trong đồ án, mỗi lệnh mang một nonce riêng; người dùng cũng có thể chủ động vô hiệu một nonce để hủy hiệu lực của lệnh mang nonce đó.

\section{Cây Merkle}
\label{apdx:merkle}

Cây Merkle~\cite{merkle1987digital} là một cấu trúc cho phép chứng minh ``một phần tử thuộc một tập hợp'' mà không cần đưa ra cả tập hợp. Từ danh sách các phần tử, ta băm từng phần tử thành các ``lá'', rồi liên tục băm từng cặp lại với nhau cho tới khi còn một giá trị duy nhất gọi là gốc (root). Để chứng minh một phần tử nằm trong cây, cần đưa ra phần tử đó cùng một số ít giá trị băm trên đường từ lá tới gốc (gọi là bằng chứng Merkle); người kiểm tra băm dần lên và so với gốc. Nhờ đó, việc lưu và ký một giá trị gốc là đủ để xác thực nhiều phần tử - đây là cách đồ án cấp quyền khớp cho từng slot xuyên chuỗi mà mỗi lệnh được ký một lần duy nhất. Một lưu ý an toàn: để chống tấn công tiền ảnh thứ hai (second-preimage), trong đó một nút trong của cây bị trình ngược lại như thể là một lá, đồ án băm mỗi lá \emph{hai lần} theo chuẩn của thư viện OpenZeppelin~\cite{openzeppelin_contracts}; nhờ đó tiền ảnh của một lá luôn dài 32 byte, khác với 64 byte của một nút trong, nên hai loại không thể bị nhầm lẫn.

\section{Hợp đồng khóa thời gian theo mã băm (HTLC)}
\label{apdx:htlc}

Hợp đồng khóa thời gian theo mã băm (Hashed Timelock Contract)~\cite{herlihy2018atomic} là một cơ chế khóa tài sản dựa trên hai điều kiện: một khóa băm và một mốc thời gian. Người ta chọn một bí mật $S$, tính khóa băm $H = \text{băm}(S)$, rồi khóa tài sản với hai điều kiện: ai đưa ra đúng $S$ được rút tài sản; nếu quá mốc thời gian mà chưa có bên nào rút, tài sản được hoàn về người gửi. Cơ chế này cho phép hai bên ở hai chuỗi khác nhau trao đổi tài sản một cách an toàn: khi một bên tiết lộ $S$ để nhận tài sản ở chuỗi này, bên kia đọc được $S$ và dùng nó để nhận tài sản ở chuỗi kia. Nhờ vậy, hai bên hoặc cùng nhận được tài sản, hoặc cùng được hoàn lại, mà không cần một bên trung gian giữ tiền.

\section{Chuẩn ký EIP-712 và Permit2}
\label{apdx:eip712}

EIP-712 là chuẩn ký dữ liệu có cấu trúc trên Ethereum. Thay vì ký một chuỗi byte khó đọc, người dùng ký một cấu trúc dữ liệu có tên trường rõ ràng, nên ví có thể hiển thị đúng nội dung sắp ký, giảm nguy cơ ký nhầm nội dung độc hại. Bản tin ký còn gắn với một ``miền'' (domain) gồm tên ứng dụng, mã chuỗi và địa chỉ hợp đồng, để một chữ ký dùng cho ứng dụng này không dùng lại được cho ứng dụng khác.

Permit2 là một hợp đồng cho phép người dùng cấp quyền chi tiêu token một lần, sau đó các hợp đồng được ủy quyền có thể kéo token tới đúng địa chỉ đích mà không cần người dùng ký lại từng lần. Đồ án dùng Permit2 để kéo tài sản của người dùng vào đúng lúc một phần lệnh được khớp, nhờ đó người dùng không phải khóa tài sản trước và không tốn phí gas khi tạo lệnh.

\section{Đấu giá kiểu Hà Lan}
\label{apdx:dutch}

Đấu giá kiểu Hà Lan (Dutch auction) là hình thức đấu giá ngược: giá bắt đầu ở mức cao rồi giảm dần theo thời gian, và người mua đầu tiên thấy mức giá đủ tốt sẽ chốt giao dịch. Khác với đấu giá tăng giá thông thường, ở đây không có việc các bên trả giá lên nhau; thị trường tự tìm ra mức giá mà tại đó có người sẵn sàng mua. Trong đồ án, giá của lệnh giảm dần theo từng block, và filler nào thấy mức giá hiện tại đủ để có lãi sẽ tham gia khớp.

\section{MEV và tấn công chèn lệnh}
\label{apdx:mev}

MEV (Maximal Extractable Value)~\cite{daian2020flashboys, qin2022quantifying} là phần giá trị mà những bên có quyền sắp xếp thứ tự giao dịch trong một block (như thợ đào hoặc người xây block) có thể trích xuất bằng cách thêm, bỏ hoặc đổi thứ tự giao dịch. Một dạng phổ biến là tấn công ``kẹp'' (sandwich): khi thấy một giao dịch mua lớn sắp được thực hiện trên một hồ thanh khoản, kẻ tấn công chèn một lệnh mua ngay trước để đẩy giá lên, để nạn nhân mua ở giá cao, rồi bán ra ngay sau đó để hưởng chênh lệch. Việc định giá trước theo thời gian như trong đấu giá kiểu Hà Lan giúp giảm bớt loại tấn công này, vì mức giá không còn được quyết định tức thời tại thời điểm giao dịch.

\section{CREATE2}
\label{apdx:create2}

CREATE2 là một opcode của EVM (được giới thiệu trong EIP-1014) cho phép triển khai hợp đồng tới một địa chỉ xác định trước mà không cần thực sự triển khai. Địa chỉ được tính từ bốn thành phần: tiền tố cố định \texttt{0xff}, địa chỉ hợp đồng triển khai, một giá trị \texttt{salt} tùy chọn, và hàm băm của bytecode hợp đồng:
\begin{multline}
    \texttt{addr} = \texttt{keccak256}\bigl(\texttt{0xff} \;\|\; \texttt{deployer} \;\|\; \texttt{salt}\\
    \|\; \texttt{keccak256(bytecode)}\bigr)[12:]
\end{multline}

Tính chất này cho phép hai bên tính toán và đồng thuận về địa chỉ hợp đồng ký quỹ \emph{trước} khi nó được triển khai. Trong đồ án, \texttt{EscrowSrcFactory} và \texttt{EscrowDstFactory} dùng CREATE2 để tạo các hợp đồng ký quỹ theo slot: salt được tính từ thông tin lệnh và chỉ số slot, đảm bảo địa chỉ hợp đồng là duy nhất và có thể xác minh độc lập bởi cả hai bên mà không cần giao tiếp trước.

\section{Các thư viện OpenZeppelin sử dụng}
\label{apdx:openzeppelin}

Đồ án sử dụng thư viện OpenZeppelin Contracts~\cite{openzeppelin_contracts} — tập hợp các thành phần hợp đồng thông minh đã được kiểm toán và được dùng rộng rãi trong thực tế. Bảng~\ref{table:oz_libs} liệt kê các thư viện được dùng trực tiếp.

\begin{table}[H]
    \centering
    \caption{Các thư viện OpenZeppelin được sử dụng trong đồ án.}
    \label{table:oz_libs}
    \begin{tabular}{|l|p{9.5cm}|}
        \hline
        \textbf{Thư viện} & \textbf{Vai trò trong đồ án} \\
        \hline
        \texttt{MerkleProof} & Xác minh bằng chứng Merkle cho từng slot xuyên chuỗi; băm lá hai lần để chống tấn công tiền ảnh thứ hai. \\
        \hline
        \texttt{ECDSA} & Phục hồi địa chỉ người ký từ chữ ký EIP-712 để xác thực chữ ký của người dùng và cosigner. \\
        \hline
        \texttt{SafeERC20} & Bọc các lần gọi \texttt{transfer}/\texttt{transferFrom} để phát hiện token trả về \texttt{false} thay vì revert, tránh mất tài sản thầm lặng. \\
        \hline
        \texttt{ReentrancyGuard} & Cung cấp modifier \texttt{nonReentrant} bảo vệ các hàm khớp lệnh và rút tài sản khỏi tấn công gọi lại. \\
        \hline
    \end{tabular}
\end{table}

\end{document}
